% ======================================================================
%  REACH Center Climate and Health Research Day
%  Poster: Climate Swings, Sticky Costs
%  Author: Ahwaz Akhtar (George Washington University, ahwaz@gwu.edu)
%  Format: 48 in x 36 in (landscape), US academic standard
%  Compile: pdflatex reach_poster.tex  (twice if labels are added later)
% ======================================================================

\documentclass[final,t]{beamer}

% --- Poster size: 48 in = 121.92 cm, 36 in = 91.44 cm ---
\usepackage[orientation=landscape,size=custom,width=121.92,height=91.44,scale=1.22]{beamerposter}

% --- Fonts & typography ---
\usepackage[T1]{fontenc}
\usepackage[utf8]{inputenc}
\usepackage{lmodern}
\usepackage{helvet}
\renewcommand{\familydefault}{\sfdefault}
\usepackage{microtype}

% --- Graphics & layout ---
\usepackage{graphicx}
\usepackage{booktabs}
\usepackage{array}
\usepackage{tabularx}
\usepackage{xcolor}
\usepackage{tcolorbox}
\tcbuselibrary{skins,breakable}
\usepackage{multicol}
\usepackage{amsmath}
\usepackage{amssymb}
\usepackage{ragged2e}

% ======================================================================
%  COLOR PALETTE  (George Washington University)
% ======================================================================
\definecolor{gwblue}{HTML}{033C5A}   % GWU primary
\definecolor{gwbuff}{HTML}{AC935E}   % GWU secondary
\definecolor{accent}{HTML}{B91C1C}   % headline accent (red)
\definecolor{softgray}{HTML}{EEEEEE}
\definecolor{darkgray}{HTML}{444444}
\definecolor{bodytext}{HTML}{222222}

% ======================================================================
%  BEAMER THEME OVERRIDES
% ======================================================================
\usetheme{default}
\setbeamercolor{background canvas}{bg=white}
\setbeamertemplate{navigation symbols}{}
\setbeamertemplate{itemize item}{\color{gwblue}\rule[0.25ex]{0.9ex}{0.9ex}}
\setbeamertemplate{itemize subitem}{\color{gwbuff}\textendash}

% ======================================================================
%  POSTER BLOCK COMMAND  (tcolorbox)
% ======================================================================
% Title sits INSIDE the box (flush) so long titles can't clip past the column edge.
\newtcolorbox{posterblock}[1]{
  enhanced,
  colback=white,
  colframe=gwblue,
  boxrule=0.8mm,
  arc=2mm,
  left=8mm,right=8mm,top=4mm,bottom=5mm,
  title={\MakeUppercase{#1}},
  fonttitle=\bfseries\large,
  colbacktitle=gwblue,
  coltitle=white,
  sharp corners=northwest,sharp corners=northeast,
  toptitle=3mm,bottomtitle=3mm,
  before upper={\vspace{1mm}},
  breakable
}

% ======================================================================
%  CALLOUT BOX  (for key findings)
% ======================================================================
\newtcolorbox{callout}{
  enhanced,
  colback=gwbuff!15,
  colframe=gwbuff,
  boxrule=0.6mm,
  arc=2mm,
  left=6mm,right=6mm,top=4mm,bottom=4mm,
}

% ======================================================================
%  DOCUMENT
% ======================================================================
\begin{document}
\begin{frame}[t,fragile]

% ----------------------------------------------------------------------
%  HEADER BAND
% ----------------------------------------------------------------------
\begin{tcolorbox}[
  enhanced,
  colback=gwblue,
  colframe=gwblue,
  arc=0mm,
  left=10mm,right=10mm,top=6mm,bottom=6mm,
  boxrule=0pt
]
\begin{columns}[c]
  \begin{column}{0.80\paperwidth}
    \centering
    \vspace{-2mm}
    {\color{white}\fontsize{82pt}{92pt}\selectfont\bfseries Climate Swings, Sticky Costs\par}
    \vspace{4mm}
    {\color{gwbuff}\fontsize{38pt}{46pt}\selectfont
    Year\mbox{-}over\mbox{-}Year Weather Volatility and the Financial Health of U.S.\ Counties, 2011--2023\par}
    \vspace{8mm}
    {\color{white}\fontsize{30pt}{36pt}\selectfont
    \textbf{Ahwaz Akhtar} \enspace\textbar\enspace George Washington University
    \enspace\textbar\enspace \texttt{ahwaz@gwu.edu}
    \enspace\textbar\enspace \textit{REACH Center Climate and Health Research Day, April 2026}\par}
    \vspace{-2mm}
  \end{column}
  \begin{column}{0.17\paperwidth}
    \centering
    \colorbox{white}{%
      \hspace{2mm}%
      \includegraphics[trim={30 80 660 50}, clip, height=6.5cm]{ocm_primary.png}%
      \hspace{2mm}%
    }
  \end{column}
\end{columns}
\end{tcolorbox}

\vspace{8mm}

% ----------------------------------------------------------------------
%  THREE-COLUMN BODY
% ----------------------------------------------------------------------
\begin{columns}[t]

% ============================================================
%  LEFT COLUMN  — Motivation / Data / Method
% ============================================================
\begin{column}{0.305\paperwidth}

\begin{posterblock}{Motivation}
\normalsize
Climate change is reshaping U.S.\ health finance not only through extreme events but through \textbf{year-to-year volatility}. Insurers reprice on new information, hospitals absorb uncompensated care, and households accumulate medical debt with a lag. Yet most economic studies estimate \emph{level} effects of heat, drought, or pollution, treating a county with persistent exposure the same as one in transition.

\vspace{2mm}
\textbf{This poster asks:} does the \emph{change} in exposure from one year to the next \textemdash{} a weather swing \textemdash{} independently burden the health-finance system, above and beyond the level?

\vspace{2mm}
If costs are driven by volatility, \textbf{adaptation planning must price in transitions, not just averages}.
\end{posterblock}

\vspace{3mm}

\begin{posterblock}{Data}
\normalsize
\begin{itemize}
\setlength\itemsep{1mm}
\item \textbf{Panel:} 3{,}155 U.S.\ counties \texttimes{} 2011--2023 (41{,}376 county-years)
\item \textbf{Climate} (NOAA National Centers for Environmental Information): divisional temperature, precipitation, Cooling Degree Days (CDD), Heating Degree Days (HDD), Palmer Drought Severity Index (PDSI); 1990--2000 baseline for county-specific z-scores
\item \textbf{Air quality} (U.S.\ Environmental Protection Agency, EPA): annual Median and Max Air Quality Index (AQI) by county
\item \textbf{Health-finance outcomes:} medical debt in collections (Urban Institute); Benchmark Silver premium (Health Insurance Exchange data, HIX Compare; rating-area \textrightarrow{} county); hospital bad debt per capita (National Academy for State Health Policy Hospital Cost Tool, NASHP HCT)
\item \textbf{Economic outcomes:} Per-Capita Personal Income (Bureau of Economic Analysis, BEA); median household income and civilian employment (U.S.\ Census American Community Survey, ACS)
\end{itemize}
\end{posterblock}

\vspace{3mm}

\begin{posterblock}{Empirical Strategy}
\normalsize
The primary specification isolates the marginal effect of a weather \emph{swing} while controlling for the prior year's level:

\vspace{1mm}
\begin{center}
\setlength{\fboxsep}{3mm}
\colorbox{softgray}{%
\parbox{0.92\linewidth}{%
\centering\large
$Y_{it} = \beta_1\,\Delta X_{it} + \beta_2\,X_{i,t-1} + \gamma\mathbf{Z}_{it} + \alpha_i + \tau_t + \varepsilon_{it}$}}
\end{center}
\vspace{1mm}

\begin{itemize}
\setlength\itemsep{1mm}
\item $\boldsymbol{\beta_1}$ \textbf{(key):} effect of a one-unit change in exposure from $t{-}1$ to $t$, holding last year's level fixed. $\boldsymbol{\beta_2}$ absorbs the residual level effect.
\item $\boldsymbol{\alpha_i,\,\tau_t}$: county FE absorb time-invariant heterogeneity (climate zone, baseline health, demographics); year FE absorb national shocks (inflation, CRA changes). $\mathbf{Z}_{it}$ controls: real PCPI, household income, employment.
\item \textbf{Clustering:} state-level primary; rating-area for premium outcomes that share rating-area pricing.
\item \textbf{Local projections} (Jordà 2005) extend to horizons $h=0,1,2,3$, tracing dynamic impulse-response without imposing lag form.
\item \textbf{Decompositions:} asymmetric splits $\Delta X$ into positive (escalation) vs negative (relief); onset/exit/persist isolates regime transitions.
\item \textbf{Exclusions:} CO 2023 dropped (HB23-1126) so a debt-reporting policy change is not read as a climate effect.
\end{itemize}

\vspace{1mm}
\textbf{896 coefficient estimates} across 5 specs \texttimes{} 7 outcomes \texttimes{} 7 exposures \texttimes{} 4 horizons.
\end{posterblock}

\end{column}

% ============================================================
%  MIDDLE COLUMN  — Finding 1 (AQI escalation + ratchet)
% ============================================================
\begin{column}{0.335\paperwidth}

\begin{posterblock}{Finding 1 \textemdash{} AQI Swings Compound Over Time}
\normalsize
A one-unit positive swing in a county's Max AQI is associated with a \textbf{rising} stream of hospital bad debt across all four horizons, with the effect \emph{growing} over time (h=0: +0.005, p=0.001; h=3: +0.016, p\textless.0001). There is no recovery \textemdash{} the cost accumulates.

\vspace{2mm}
\begin{center}
\includegraphics[width=0.78\linewidth]{plots/F1_aqi_compound.png}
\end{center}

\begin{callout}
\normalsize\centering
\textbf{A worsening air-quality year today adds to hospital bad debt \emph{three years later} \textemdash{} after claims mature, patients default, and losses are recognized.}
\end{callout}
\end{posterblock}

\vspace{4mm}

\begin{posterblock}{Finding 2 \textemdash{} The Ratchet: Relief Does Not Reverse}
\normalsize
Splitting the swing by direction uncovers a counterintuitive pattern. \textbf{A year in which a county's air quality \emph{improves} still raises hospital bad debt by +\$0.010 per capita} (p\textless.0001). A year in which air quality \emph{worsens} shows no significant contemporaneous effect. Bad debt climbs regardless of this year's trend \textemdash{} prior pollution exposure keeps converting into uncompensated care.

\vspace{2mm}
\begin{center}
\includegraphics[width=0.78\linewidth]{plots/F2_aqi_ratchet.png}
\end{center}

\begin{callout}
\normalsize\centering
\textbf{The cost is a ratchet, not a thermostat.} Improving air quality does not unwind hospital bad debt within three years \textemdash{} prior exposure continues to mature into new uncompensated-care losses.
\end{callout}

\end{posterblock}

\end{column}

% ============================================================
%  RIGHT COLUMN  — Premiums / Income / Takeaways
% ============================================================
\begin{column}{0.305\paperwidth}

\begin{posterblock}{Finding 3 \textemdash{} Cold-Shock \emph{Onset} Prices Premiums}
\normalsize
The first year a county \emph{enters} a high-HDD regime: \textbf{+\$49.46 on the Benchmark Silver premium} (p=0.003). Exit and persistence are null \textemdash{} insurers price at the transition, then hold.

\vspace{1mm}
\begin{center}
\includegraphics[width=0.72\linewidth]{plots/F3_hdd_onset.png}
\end{center}
\end{posterblock}

\vspace{3mm}

\begin{posterblock}{Finding 4 \textemdash{} Drought Swings Suppress Income}
\normalsize
A worsening PDSI swing cuts per-capita personal income for two years (h=1: $-$\$226, p\textless.0001; h=2: $-$\$144, p=0.002); concentrated in \textbf{drought-worsening swings} ($-$\$207, p=0.001), with employment loss ($-$142 jobs, p=0.006).

\vspace{1mm}
\begin{center}
\includegraphics[width=0.72\linewidth]{plots/F4_pdsi_income.png}
\end{center}
\end{posterblock}

\vspace{3mm}

\begin{posterblock}{Takeaways}
\normalsize
\begin{itemize}
\setlength\itemsep{2mm}
\item \textbf{Volatility is its own exposure.} The swing channel is significant where the level channel was not (AQI $\to$ hospital bad debt).
\item \textbf{Costs ratchet, not reverse.} Pollution relief does not refund hospital bad debt within three years; insurers reprice on shock \emph{onset}, not exit.
\item \textbf{Drought volatility hits the macro base.} Income and employment contract for multiple years, feeding back into medical-debt formation.
\end{itemize}
\end{posterblock}

\end{column}
\end{columns}

% ----------------------------------------------------------------------
%  ACKNOWLEDGMENTS BAND
% ----------------------------------------------------------------------
\vspace{5mm}
\begin{tcolorbox}[
  enhanced,
  colback=gwblue,
  colframe=gwblue,
  arc=0mm,
  left=10mm,right=10mm,top=3mm,bottom=3mm,
  boxrule=0pt
]
\centering\color{white}\normalsize
\textbf{Acknowledgments.}\enspace Thanks to faculty \textbf{Avi Dor}, \textbf{Susan Anenberg}, \textbf{Eric Luo}, and \textbf{Cindy Liu} for guidance on specification, data, and interpretation.
\end{tcolorbox}

\end{frame}
\end{document}
